% FRAGMENT A15 (drafter A: class model and main results)
% REPLACES: main.tex lines 1925–1948 ("\subsubsection{Nonemptiness for every value" to "… not a statement about effect size at every value ratio.")
%   The blank line 1949 and the heading of Appendix A.5 at line 1950 stay.
% CONTENT: (1) nonemptiness of the Proposition 2 region for every ranked pair of classes,
%   every prior, noise scale, and low-cost probability, with part (iii) and the ownership
%   condition; (2) a short subsection on more than two classes.
% LABELS KEPT: nonemptiness-for-every-value-gap (title changes), eq:paper-a-25.
% LABELS REMOVED: none.
% LABELS ADDED: more-than-two-classes (new subsubsection; nothing references it yet).
% NEW SYMBOLS: \Gamma(x)=\int_p^x D(u)du and its limit \Gamma=E_H max{θ,p}−E_L max{θ,p};
%   u_0 (lowest value above which the classes differ); F_r (uniform law on [0,r]);
%   W_H(r) (as W_H(F) in Proposition 2, at F=F_r). Many-class subsection: \pi_j (class
%   priors), T_X (posterior mean of class proceeds), \bar e(q) (expected entry rate at order
%   q), \bar\pi (prior mass of the high group), \mathcal M_\pi (feasible posterior set),
%   \gamma (separation of the group proceeds from the feasible posterior means).
%   The source memos use other letters (a reserved Greek letter for the group posterior,
%   \alpha for the group prior, D_0 and d_1 for the ranges); those letters collide with
%   \alpha_H, D(u), and the brief's reserved symbol, so this fragment does not use them.
% NEW CITATIONS: none.  NUM PLACEHOLDERS: none.
% SOURCES: Q1_N_main.md Claim 11 (nonemptiness); Q1_N_extensions.md Claims 7, 8, 9
%   (many classes); SYNTHESIS.md A.1 rows 4, 25.

\subsubsection{Nonemptiness for every ranked pair of
classes}\label{nonemptiness-for-every-value-gap}

The region in \hyperref[result-proposition-2]{Proposition 2} needs no
large class gap and no special prior. Fix classes under
\hyperref[sec-values-preparation-and-the-sale]{Section 2.2}. Set
\(\Gamma(x)=\int_p^xD(u)\,du\). It is nondecreasing, with limit
\(\Gamma=\mathbb E_H[\max\{\theta,p\}]-\mathbb E_L[\max\{\theta,p\}]\).
Take classes that differ above the reserve, \(\Gamma>0\). Fix
\(\pi\in(0,1)\), \(b>1\), and \(\rho\in(0,1)\). Let
\(u_0=\inf\{x\ge p:\Gamma(x)>0\}\), the lowest value above which the
classes differ. In the point-mass benchmark \(u_0=\ell\) and
\(\Gamma=h-\ell\). Let \(F_r\) be the uniform distribution on \([0,r]\)
with \(r\ge p\). The kernel \(\xi\) in the proof of
\hyperref[result-proposition-1]{Proposition 1} equals \(\Gamma(R)\) for
\(R\ge p\) and, under equal admissibility, zero for \(R<p\), so
\(\Delta_T(r)=r^{-1}\int_p^r\Gamma(u)\,du\). Hence \(\Delta_T(r)=0\) for
\(p\le r\le u_0\), and for \(r>u_0\) the spread is continuous and
strictly increasing, with derivative
\(r^{-2}[r\Gamma(r)-\int_p^r\Gamma(u)\,du]\ge p\Gamma(r)/r^2>0\). Since
\(\Gamma(x)\le\Gamma\), \(\Delta_T(r)\le(1-p/r)\Gamma\), and the split
in the proof of Proposition 1 gives
\(g_H-g_L=\Gamma-\Delta_T(r)\ge\Gamma p/r>0\). Dominated convergence,
with the finite class means, makes \(r\mapsto g_j(r)\) and
\(r\mapsto W_H(r)=\mathbb E_H[F_r(\theta)\mathbf1\{\theta\ge p\}]\)
continuous and gives \(g_j(r)\to0\) as \(r\to\infty\). Also
\(W_H(u_0)\ge(p/u_0)\Pr_H(\theta\ge p)>0\), because \(\Gamma>0\) gives
the high class mass above the reserve. At the limit \(r=u_0\),

\[
g_H-g_L=\Gamma,\qquad
B_{u_0}(M(\pi))-B_{u_0}(\pi)=[M(\pi)-\pi]\Gamma>0.
\tag{A.25}\label{eq:paper-a-25}
\]

Choose \(r_1>u_0\) close enough to \(u_0\) that
\(B_{r_1}(M(\pi))>B_{u_0}(\pi)\) and
\((1+\rho)W_H(r_1)>2\rho W_H(u_0)\). Since \(B_r(\pi)\) and \(W_H(r)\)
decrease in \(r\), every \(r_0\in(u_0,r_1)\) satisfies
\(B_{r_0}(\pi)<B_{r_1}(M(\pi))\) and
\((1+\rho)W_H(r_1)>2\rho W_H(r_0)>0\). Choose \(r_0\) close enough to
\(u_0\) that \(\Delta_T(r_0)<(1-1/b)\rho\kappa(\pi)\Delta_T(r_1)\),
which is possible because \(\Delta_T(r)\to0\) as \(r\downarrow u_0\).
Select \(k\) strictly inside the window of \eqref{eq:paper-a3} and
\(c_H\) strictly inside the window of \eqref{eq:paper-a2}; then
\(c_H>B_{r_0}(\pi)>0\). Since \(B_r(M(\pi))\to0\) as \(r\to\infty\),
some \(r_2>r_1\) has \(B_{r_2}(M(\pi))<c_H\), and
\(\Delta_T(r_2)\ge\Delta_T(r_1)\) gives the trading bound of part
(iii). All inequalities have strict slack and persist under small
perturbations of \((r_0,r_1,r_2,c_H,k)\) and of \((\pi,b,\rho)\), with
the low cost fixed at zero. This is a mathematical nonemptiness
construction, not a statement about effect size for every pair of
classes.

\subsubsection{More than two classes}\label{more-than-two-classes}

Price sufficiency does not need two classes, and the main comparison
extends to two groups of classes. Let the class take finitely many
values \(j\) with prior probabilities \(\pi_j>0\), class value
distributions \(G_j\) with finite means, and class objects \(t_j,g_j\)
as in \hyperref[sec-equilibrium]{Section 2.4}; the investor observes
the class. The low cost is zero, so entry at every price is at least
\(\rho\), and competitive pricing gives \(P=t_0+e(P)(T_X-t_0)\), where
\(T_X\) is the posterior mean of the class proceeds at flow \(X\). The
price therefore reveals \(T_X\).

\emph{Price sufficiency with many classes (analytical).} In every
equilibrium the price reveals the whole posterior over classes, with no
restriction on how the class profits \(g_j\) relate to the class
proceeds \(t_j\). Fix a candidate schedule, and let
\(\bar e(q)=\int e(P(x))f(x-q)\,dx\) be the expected entry rate when the
order is \(q\). A class-\(j\) investor who deviates to \(q\) earns
\(t_jq\bar e(q)\) minus a term that does not depend on the class. The
translation bound gives
\([q\bar e(q)]'\ge(1-|q|/b)\bar e(q)\ge(1-1/b)\rho>0\), so
\(q\bar e(q)\) is strictly increasing. Adding the optimality conditions
of two classes with \(t_j<t_l\) gives
\((t_l-t_j)[q_l\bar e(q_l)-q_j\bar e(q_j)]\ge0\), so the class with
lower proceeds never chooses a larger order, also within mixtures. For
flows \(x_1<x_2\) the likelihood ratio \(f(x_2-q)/f(x_1-q)\) is
nondecreasing in \(q\), and

\[
T_X(x_2)-T_X(x_1)
=\frac{\operatorname{Cov}_{x_1}\!\left[t_j,\,f(x_2-q)/f(x_1-q)\right]}
{\mathbb E_{x_1}\!\left[f(x_2-q)/f(x_1-q)\right]}\ge0,
\]

where the moments are over class and order under the posterior at
\(x_1\). The covariance is half the expectation of the product of the
difference in class proceeds and the difference in likelihood ratios
over two independent draws, and every such product is nonnegative. If
the class proceeds are not all equal, equality forces the likelihood
ratio to be constant on the joint support of class and order, and then
the whole posterior is the same at \(x_1\) and \(x_2\). The posterior is
therefore a measurable function of \(T_X\), hence of the price. If all
class proceeds are equal, no order earns a residual and the posterior
stays at the prior.

\emph{Two groups of classes (analytical).} Split the classes into a low
group and a high group, with prior mass \(\bar\pi\) on the high group.
Let \(\mathcal M_\pi\) be the set of posteriors whose likelihood ratios
relative to the prior differ by at most the factor \(e^{2/b}\) across
classes; every flow and price posterior lies in \(\mathcal M_\pi\). In
the strong economy let \(\gamma\) be the smallest distance between a
class's proceeds and the range \(\{\mu\cdot t:\mu\in\mathcal M_\pi\}\)
of feasible posterior means, with the low group below that range and
the high group above it. Suppose \(\gamma>0\) and
\(\max_jt_j-\min_jt_j<k<(1-1/b)\rho\gamma\), with the first difference
taken in the weak economy; gross profit at the prior is below \(c_H\) in
the weak economy; and in the strong economy the high group has the
higher mean profit, with gross profit above \(c_H\) at the group
posterior \(M(\bar\pi)\). Then the weak economy has unique zero orders
and entry \(\rho\), and the strong economy has unique orders \(-1\) in
the low group and \(+1\) in the high group, an informative price, and
entry above \(\rho\). The proof repeats Steps 3 to 6 of the proof of
\hyperref[result-proposition-2]{Proposition 2}. In the weak economy
every residual is at most \(\max_jt_j-\min_jt_j\) in absolute value. In
the strong economy every residual has the sign of its group and
magnitude at least \(\rho\gamma\). Under the forced orders, flow
carries information only about the group, so expected proceeds and
gross profit are increasing affine functions of the group posterior,
which lies in \([m(\bar\pi),M(\bar\pi)]\). With two classes
\(\gamma=\kappa(\pi)\Delta_T\), and the conditions reduce to
\eqref{eq:paper-a2} and \eqref{eq:paper-a3}. Ownership by the high group
needs a separate comparison, as in part (ii) of Proposition 2. The
argument does not cover a class whose proceeds lie inside the range of
feasible posterior means: its residual can change sign with flow, and
its equilibrium orders remain open.
